\documentclass[10pt,a4paper]{amsart}
\usepackage{amsmath,amssymb,amsthm,graphicx}
% Retained minimal preamble for the Stacks Project chapter flat.tex (More on Flatness).
% Verbatim excerpt of the chapter preamble: Xy-pic setup (native preamble lines 14-18) and the native
% theorem declarations with the mathematical macros (native preamble lines 154-217). Copied mechanically;
% nothing added, removed or redefined. The wrapper class and the amsmath/amssymb/amsthm/graphicx base
% packages are supplied by the assembler, so the class/externaldocument part of the native preamble is omitted.
% For commutative diagrams we use Xy-pic
\usepackage[all]{xy}

% We use 2cell for 2-commutative diagrams.
\xyoption{2cell}
% Theorem environments.
%
\theoremstyle{plain}
\newtheorem{theorem}[subsection]{Theorem}
\newtheorem{proposition}[subsection]{Proposition}
\newtheorem{lemma}[subsection]{Lemma}

\theoremstyle{definition}
\newtheorem{definition}[subsection]{Definition}
\newtheorem{example}[subsection]{Example}
\newtheorem{exercise}[subsection]{Exercise}
\newtheorem{situation}[subsection]{Situation}

\theoremstyle{remark}
\newtheorem{remark}[subsection]{Remark}
\newtheorem{remarks}[subsection]{Remarks}

\numberwithin{equation}{subsection}

% Macros
%
\def\lim{\mathop{\mathrm{lim}}\nolimits}
\def\colim{\mathop{\mathrm{colim}}\nolimits}
\def\Spec{\mathop{\mathrm{Spec}}}
\def\Hom{\mathop{\mathrm{Hom}}\nolimits}
\def\Ext{\mathop{\mathrm{Ext}}\nolimits}
\def\SheafHom{\mathop{\mathcal{H}\!\mathit{om}}\nolimits}
\def\SheafExt{\mathop{\mathcal{E}\!\mathit{xt}}\nolimits}
\def\Sch{\mathit{Sch}}
\def\Mor{\mathop{\mathrm{Mor}}\nolimits}
\def\Ob{\mathop{\mathrm{Ob}}\nolimits}
\def\Sh{\mathop{\mathit{Sh}}\nolimits}
\def\NL{\mathop{N\!L}\nolimits}
\def\CH{\mathop{\mathrm{CH}}\nolimits}
\def\proetale{{pro\text{-}\acute{e}tale}}
\def\etale{{\acute{e}tale}}
\def\QCoh{\mathit{QCoh}}
\def\Ker{\mathop{\mathrm{Ker}}}
\def\Im{\mathop{\mathrm{Im}}}
\def\Coker{\mathop{\mathrm{Coker}}}
\def\Coim{\mathop{\mathrm{Coim}}}

% Boxtimes
%
\DeclareMathSymbol{\boxtimes}{\mathbin}{AMSa}{"02}

%
% Macros for moduli stacks/spaces
%
\def\QCohstack{\mathcal{QC}\!\mathit{oh}}
\def\Cohstack{\mathcal{C}\!\mathit{oh}}
\def\Spacesstack{\mathcal{S}\!\mathit{paces}}
\def\Quotfunctor{\mathrm{Quot}}
\def\Hilbfunctor{\mathrm{Hilb}}
\def\Curvesstack{\mathcal{C}\!\mathit{urves}}
\def\Polarizedstack{\mathcal{P}\!\mathit{olarized}}
\def\Complexesstack{\mathcal{C}\!\mathit{omplexes}}
% \Pic is the operator that assigns to X its picard group, usage \Pic(X)
% \Picardstack_{X/B} denotes the Picard stack of X over B
% \Picardfunctor_{X/B} denotes the Picard functor of X over B
\def\Pic{\mathop{\mathrm{Pic}}\nolimits}
\def\Picardstack{\mathcal{P}\!\mathit{ic}}
\def\Picardfunctor{\mathrm{Pic}}
\def\Deformationcategory{\mathcal{D}\!\mathit{ef}}

% --- appended 2026-10-12 (Chief): the project own definitions of the reference/slogan/history
% display environments, copied mechanically from the shared Stacks preamble lines 8-12
% (\usepackage{verbatim} plus the three \newenvironment lines, which the chapter preamble keeps
% verbatim from the project preamble). Nothing added, removed or redefined: these environments
% carry no mathematical content and the project itself typesets them as comments. Without these
% three lines a counted statement whose own text uses them fails with Environment slogan undefined.
\usepackage{verbatim}
\newenvironment{reference}{\comment}{\endcomment}
\newenvironment{slogan}{\comment}{\endcomment}
\newenvironment{history}{\comment}{\endcomment}
\usepackage{cleveref}
\crefname{originallemma}{original lemma}{original lemmas}
\Crefname{originallemma}{Original lemma}{Original lemmas}
\makeatletter
\newlabel{algebra-lemma-ass-weakly-ass}{{058A}{source}{}{sourcepointer-1}{}}
\newlabel{algebra-lemma-ass-weakly-ass@cref}{{[originallemma][0][]058A}{[1][0][]source}{}{}{}}
\makeatother
\crefname{originallemma}{original lemma}{original lemmas}
\Crefname{originallemma}{Original lemma}{Original lemmas}
\makeatletter
\newlabel{algebra-lemma-ass-zero-divisors}{{00LD}{source}{}{sourcepointer-2}{}}
\newlabel{algebra-lemma-ass-zero-divisors@cref}{{[originallemma][0][]00LD}{[1][0][]source}{}{}{}}
\makeatother
\crefname{originallemma}{original lemma}{original lemmas}
\Crefname{originallemma}{Original lemma}{Original lemmas}
\makeatletter
\newlabel{algebra-lemma-etale-at-prime}{{00U4}{source}{}{sourcepointer-3}{}}
\newlabel{algebra-lemma-etale-at-prime@cref}{{[originallemma][0][]00U4}{[1][0][]source}{}{}{}}
\makeatother
\crefname{originallemma}{original lemma}{original lemmas}
\Crefname{originallemma}{Original lemma}{Original lemmas}
\makeatletter
\newlabel{algebra-lemma-finite-ass}{{00LC}{source}{}{sourcepointer-4}{}}
\newlabel{algebra-lemma-finite-ass@cref}{{[originallemma][0][]00LC}{[1][0][]source}{}{}{}}
\makeatother
\crefname{originallemma}{original lemma}{original lemmas}
\Crefname{originallemma}{Original lemma}{Original lemmas}
\makeatletter
\newlabel{algebra-lemma-silly}{{00DS}{source}{}{sourcepointer-5}{}}
\newlabel{algebra-lemma-silly@cref}{{[originallemma][0][]00DS}{[1][0][]source}{}{}{}}
\makeatother
\crefname{originallemma}{original lemma}{original lemmas}
\Crefname{originallemma}{Original lemma}{Original lemmas}
\makeatletter
\newlabel{algebra-lemma-weakly-ass}{{0548}{source}{}{sourcepointer-6}{}}
\newlabel{algebra-lemma-weakly-ass@cref}{{[originallemma][0][]0548}{[1][0][]source}{}{}{}}
\makeatother
\crefname{originallemma}{original lemma}{original lemmas}
\Crefname{originallemma}{Original lemma}{Original lemmas}
\makeatletter
\newlabel{algebra-lemma-weakly-ass-local}{{0566}{source}{}{sourcepointer-7}{}}
\newlabel{algebra-lemma-weakly-ass-local@cref}{{[originallemma][0][]0566}{[1][0][]source}{}{}{}}
\makeatother
\crefname{originallemma}{original lemma}{original lemmas}
\Crefname{originallemma}{Original lemma}{Original lemmas}
\makeatletter
\newlabel{algebra-theorem-projective-free-over-local-ring}{{0593}{source}{}{sourcepointer-8}{}}
\newlabel{algebra-theorem-projective-free-over-local-ring@cref}{{[originallemma][0][]0593}{[1][0][]source}{}{}{}}
\makeatother
\crefname{originallemma}{original lemma}{original lemmas}
\Crefname{originallemma}{Original lemma}{Original lemmas}
\makeatletter
\newlabel{lemma-etale-weak-assassin-up-down}{{05FP}{source}{}{sourcepointer-9}{}}
\newlabel{lemma-etale-weak-assassin-up-down@cref}{{[originallemma][0][]05FP}{[1][0][]source}{}{}{}}
\makeatother
\crefname{originallemma}{original lemma}{original lemmas}
\Crefname{originallemma}{Original lemma}{Original lemmas}
\makeatletter
\newlabel{lemma-invert-universally-injective}{{05FR}{source}{}{sourcepointer-10}{}}
\newlabel{lemma-invert-universally-injective@cref}{{[originallemma][0][]05FR}{[1][0][]source}{}{}{}}
\makeatother
\crefname{originallemma}{original lemma}{original lemmas}
\Crefname{originallemma}{Original lemma}{Original lemmas}
\makeatletter
\newlabel{lemma-weak-bourbaki-pre-pre}{{0CU6}{source}{}{sourcepointer-11}{}}
\newlabel{lemma-weak-bourbaki-pre-pre@cref}{{[originallemma][0][]0CU6}{[1][0][]source}{}{}{}}
\makeatother
\crefname{originallemma}{original lemma}{original lemmas}
\Crefname{originallemma}{Original lemma}{Original lemmas}
\makeatletter
\newlabel{more-algebra-definition-auto-ass}{{05GM}{source}{}{sourcepointer-12}{}}
\newlabel{more-algebra-definition-auto-ass@cref}{{[originallemma][0][]05GM}{[1][0][]source}{}{}{}}
\makeatother
\crefname{originallemma}{original lemma}{original lemmas}
\Crefname{originallemma}{Original lemma}{Original lemmas}
\makeatletter
\newlabel{more-algebra-lemma-P-fPD-zero}{{05GQ}{source}{}{sourcepointer-13}{}}
\newlabel{more-algebra-lemma-P-fPD-zero@cref}{{[originallemma][0][]05GQ}{[1][0][]source}{}{}{}}
\makeatother
\crefname{originallemma}{original lemma}{original lemmas}
\Crefname{originallemma}{Original lemma}{Original lemmas}
\makeatletter
\newlabel{morphisms-lemma-etale-over-field}{{02GL}{source}{}{sourcepointer-14}{}}
\newlabel{morphisms-lemma-etale-over-field@cref}{{[originallemma][0][]02GL}{[1][0][]source}{}{}{}}
\makeatother
\crefname{originallemma}{original lemma}{original lemmas}
\Crefname{originallemma}{Original lemma}{Original lemmas}
\makeatletter
\newlabel{proposition-finite-presentation-flat-at-point}{{05ID}{source}{}{sourcepointer-15}{}}
\newlabel{proposition-finite-presentation-flat-at-point@cref}{{[originallemma][0][]05ID}{[1][0][]source}{}{}{}}
\makeatother
\title{Weakly associated primes of a tensor product with a flat finitely presented module over a finitely presented ring map}
\author{The Stacks Project Authors}
\date{stacks-project:flat.tex:local-snapshot-20261011; GNU FDL 1.2 or later}
\begin{document}
\maketitle
\noindent\textit{Editorial note.} All delivered mathematical text is the author's own, copied from the pinned Stacks Project chapter file flat.tex; no statement or proof step was rewritten, reordered or removed.\par\medskip
\noindent\textit{Editorial note.} Preamble provenance: the chapter's own retained minimal preamble (sources/stacks-flat/retained-preamble.tex, lines 1-85) is a verbatim excerpt of the native Xy-pic setup and of the native theorem declarations and mathematical macros, followed by the project's own reference/slogan/history display-environment definitions copied mechanically from the shared preamble; the wrapper class and the amsmath/amssymb/amsthm/graphicx base packages are supplied by the assembler, and extra\_packages is left empty because everything needed is already in the retained preamble. No macro was invented or redefined.\par\medskip
\noindent\textit{Editorial note.} Same-chapter pointers are located at their \textbackslash{}label lines inside the pinned chapter file; cross-chapter pointers are located in the project's own tag list staged in the snapshot (sources/stacks-chow/tags.txt), at the verbatim lines carrying their canonical tags.\par\medskip
\noindent\textit{Allowed standard prerequisites (editorial).}\begin{itemize}
\item The preliminary choice of the element g and the etale local model attached to a flat finitely presented module. Same chapter flat.tex, labels lemma-weak-bourbaki-pre-pre (invoked at proof line 3771) and proposition-finite-presentation-flat-at-point (invoked at proof line 3772). Chooses g and produces, after an etale base change R -> R' and S -> S', a situation where N tensor\_S S' is projective over R'; the key observation (equation-key-observation) is then deduced inside the counted span. The proposition is itself a separate candidate in the same pool and is consumed here only as a pointer.
\item Weakly associated primes under etale base change and universality of injectivity. Same chapter flat.tex, labels lemma-etale-weak-assassin-up-down (proof line 3799) and lemma-invert-universally-injective (proof lines 3848 and 3878). The first transports WeakAss, Ass and the weakly associated primes of the fibre through the etale maps, which is the bookkeeping of the first part of the proof; the second converts the non membership of a chosen element in the relevant primes into universal injectivity of multiplication by it, which is used twice.
\item Elementary facts on associated and weakly associated primes, projectivity over a local ring, and the auto-associated criterion for projective modules. Algebra, labels algebra-lemma-ass-weakly-ass (proof line 3809), algebra-lemma-weakly-ass-local (proof line 3832), algebra-lemma-ass-zero-divisors, algebra-lemma-finite-ass and algebra-lemma-silly (proof lines 3840-3842), algebra-lemma-etale-at-prime (proof line 3889), algebra-theorem-projective-free-over-local-ring (proof line 3899) and algebra-lemma-weakly-ass (proof line 3907); Morphisms, label morphisms-lemma-etale-over-field (proof line 3923); More on Algebra, labels more-algebra-definition-auto-ass (proof line 3958) and more-algebra-lemma-P-fPD-zero (proof line 3962). These supply the local characterisation of weakly associated primes, the existence of an element of a prime ideal that is not a zerodivisor on a given module, discreteness of the fibres of an etale morphism, freeness of a projective module over a local ring, and the criterion that a non injective map of projective modules over an auto-associated ring can be detected after tensoring with the residue field.
\end{itemize}\medskip
\ifdefined\hypertarget\hypertarget{sourcepointer-1}{}\fi
\noindent\textit{Source reference.} Original lemma 058A: sources/stacks-chow/tags.txt, source lines 6773--6773. Navigation pointer for \textbackslash{}ref\{algebra-lemma-ass-weakly-ass\} used inside the counted span (line 3809); the argument that consumes it is retained in the span.\par\medskip
\ifdefined\hypertarget\hypertarget{sourcepointer-2}{}\fi
\noindent\textit{Source reference.} Original lemma 00LD: sources/stacks-chow/tags.txt, source lines 783--783. Navigation pointer for \textbackslash{}ref\{algebra-lemma-ass-zero-divisors\} used inside the counted span (line 3840); the argument that consumes it is retained in the span.\par\medskip
\ifdefined\hypertarget\hypertarget{sourcepointer-3}{}\fi
\noindent\textit{Source reference.} Original lemma 00U4: sources/stacks-chow/tags.txt, source lines 1098--1098. Navigation pointer for \textbackslash{}ref\{algebra-lemma-etale-at-prime\} used inside the counted span (line 3889); the argument that consumes it is retained in the span.\par\medskip
\ifdefined\hypertarget\hypertarget{sourcepointer-4}{}\fi
\noindent\textit{Source reference.} Original lemma 00LC: sources/stacks-chow/tags.txt, source lines 782--782. Navigation pointer for \textbackslash{}ref\{algebra-lemma-finite-ass\} used inside the counted span (line 3841); the argument that consumes it is retained in the span.\par\medskip
\ifdefined\hypertarget\hypertarget{sourcepointer-5}{}\fi
\noindent\textit{Source reference.} Original lemma 00DS: sources/stacks-chow/tags.txt, source lines 508--508. Navigation pointer for \textbackslash{}ref\{algebra-lemma-silly\} used inside the counted span (line 3842); the argument that consumes it is retained in the span.\par\medskip
\ifdefined\hypertarget\hypertarget{sourcepointer-6}{}\fi
\noindent\textit{Source reference.} Original lemma 0548: sources/stacks-chow/tags.txt, source lines 6631--6631. Navigation pointer for \textbackslash{}ref\{algebra-lemma-weakly-ass\} used inside the counted span (line 3907); the argument that consumes it is retained in the span.\par\medskip
\ifdefined\hypertarget\hypertarget{sourcepointer-7}{}\fi
\noindent\textit{Source reference.} Original lemma 0566: sources/stacks-chow/tags.txt, source lines 6699--6699. Navigation pointer for \textbackslash{}ref\{algebra-lemma-weakly-ass-local\} used inside the counted span (line 3832); the argument that consumes it is retained in the span.\par\medskip
\ifdefined\hypertarget\hypertarget{sourcepointer-8}{}\fi
\noindent\textit{Source reference.} Original lemma 0593: sources/stacks-chow/tags.txt, source lines 6801--6801. Navigation pointer for \textbackslash{}ref\{algebra-theorem-projective-free-over-local-ring\} used inside the counted span (line 3899); the argument that consumes it is retained in the span.\par\medskip
\ifdefined\hypertarget\hypertarget{sourcepointer-9}{}\fi
\noindent\textit{Source reference.} Original lemma 05FP: sources/stacks-flat/pinned-original-source.tex, source lines 353--353. Navigation pointer for \textbackslash{}ref\{lemma-etale-weak-assassin-up-down\} used inside the counted span (line 3799); the argument that consumes it is retained in the span.\par\medskip
\ifdefined\hypertarget\hypertarget{sourcepointer-10}{}\fi
\noindent\textit{Source reference.} Original lemma 05FR: sources/stacks-flat/pinned-original-source.tex, source lines 1976--1976. Navigation pointer for \textbackslash{}ref\{lemma-invert-universally-injective\} used inside the counted span (line 3848); the argument that consumes it is retained in the span.\par\medskip
\ifdefined\hypertarget\hypertarget{sourcepointer-11}{}\fi
\noindent\textit{Source reference.} Original lemma 0CU6: sources/stacks-flat/pinned-original-source.tex, source lines 3705--3705. Navigation pointer for \textbackslash{}ref\{lemma-weak-bourbaki-pre-pre\} used inside the counted span (line 3771); the argument that consumes it is retained in the span.\par\medskip
\ifdefined\hypertarget\hypertarget{sourcepointer-12}{}\fi
\noindent\textit{Source reference.} Original lemma 05GM: sources/stacks-chow/tags.txt, source lines 7066--7066. Navigation pointer for \textbackslash{}ref\{more-algebra-definition-auto-ass\} used inside the counted span (line 3958); the argument that consumes it is retained in the span.\par\medskip
\ifdefined\hypertarget\hypertarget{sourcepointer-13}{}\fi
\noindent\textit{Source reference.} Original lemma 05GQ: sources/stacks-chow/tags.txt, source lines 7069--7069. Navigation pointer for \textbackslash{}ref\{more-algebra-lemma-P-fPD-zero\} used inside the counted span (line 3962); the argument that consumes it is retained in the span.\par\medskip
\ifdefined\hypertarget\hypertarget{sourcepointer-14}{}\fi
\noindent\textit{Source reference.} Original lemma 02GL: sources/stacks-chow/tags.txt, source lines 3232--3232. Navigation pointer for \textbackslash{}ref\{morphisms-lemma-etale-over-field\} used inside the counted span (line 3923); the argument that consumes it is retained in the span.\par\medskip
\ifdefined\hypertarget\hypertarget{sourcepointer-15}{}\fi
\noindent\textit{Source reference.} Original lemma 05ID: sources/stacks-flat/pinned-original-source.tex, source lines 3287--3287. Navigation pointer for \textbackslash{}ref\{proposition-finite-presentation-flat-at-point\} used inside the counted span (line 3772); the argument that consumes it is retained in the span.\par\medskip
\begin{lemma}
\label{lemma-weak-bourbaki-pre}
Let $R \to S$ be a ring map of finite presentation.
Let $N$ be a finitely presented $S$-module
which is flat as an $R$-module. Let $M$ be an $R$-module.
Let $\mathfrak q$ be a prime of $S$ lying over $\mathfrak p \subset R$.
Then
$$
\mathfrak q \in \text{WeakAss}_S(M \otimes_R N)
\Leftrightarrow
\Big(
\mathfrak p \in \text{WeakAss}_R(M)
\text{ and }
\overline{\mathfrak q} \in \text{Ass}_{\overline{S}}(\overline{N})
\Big)
$$
Here $\overline{S} = S \otimes_R \kappa(\mathfrak p)$,
$\overline{\mathfrak q} = \mathfrak q \overline{S}$, and
$\overline{N} = N \otimes_R \kappa(\mathfrak p)$.
\end{lemma}

\begin{proof}
Pick $g \in S$ as in Lemma \ref{lemma-weak-bourbaki-pre-pre}.
Apply Proposition \ref{proposition-finite-presentation-flat-at-point}
to the morphism of schemes $\Spec(S_g) \to \Spec(R)$, the quasi-coherent
module associated to $N_g$, and the points
corresponding to the primes $\mathfrak qS_g$ and $\mathfrak p$. Translating
into algebra we obtain a commutative diagram of rings
$$
\xymatrix{
S \ar[r] & S_g \ar[r] & S' \\
& R \ar[lu] \ar[u] \ar[r] & R' \ar[u]
}
\quad\quad
\xymatrix{
\mathfrak q \ar@{-}[r] \ar@{-}[rd] &
\mathfrak qS_g \ar@{-}[d] \ar@{-}[r] & \mathfrak q' \ar@{-}[d] \\
& \mathfrak p \ar@{-}[r] & \mathfrak p'
}
$$
endowed with primes as shown, the horizontal arrows are \'etale,
and $N \otimes_S S'$ is projective as an $R'$-module. Set
$N' = N \otimes_S S'$, $M' = M \otimes_R R'$,
$\overline{S}' = S' \otimes_{R'} \kappa(\mathfrak q')$,
$\overline{\mathfrak q}' = \mathfrak q' \overline{S}'$,
and
$$
\overline{N}' = N' \otimes_{R'} \kappa(\mathfrak p') =
\overline{N} \otimes_{\overline{S}} \overline{S}'
$$
By Lemma \ref{lemma-etale-weak-assassin-up-down} we have
\begin{align*}
\text{WeakAss}_{S'}(M' \otimes_{R'} N') & =
(\Spec(S') \to \Spec(S))^{-1}\text{WeakAss}_S(M \otimes_R N) \\
\text{WeakAss}_{R'}(M') & =
(\Spec(R') \to \Spec(R))^{-1}\text{WeakAss}_R(M) \\
\text{Ass}_{\overline{S}'}(\overline{N}') & =
(\Spec(\overline{S}') \to \Spec(\overline{S}))^{-1}
\text{Ass}_{\overline{S}}(\overline{N})
\end{align*}
Use Algebra, Lemma \ref{algebra-lemma-ass-weakly-ass}
for $\overline{N}$ and $\overline{N}'$. In particular we have
\begin{align*}
\mathfrak q \in \text{WeakAss}_S(M \otimes_R N)
& \Leftrightarrow
\mathfrak q' \in \text{WeakAss}_{S'}(M' \otimes_{R'} N') \\
\mathfrak p \in \text{WeakAss}_R(M)
& \Leftrightarrow
\mathfrak p' \in \text{WeakAss}_{R'}(M') \\
\overline{\mathfrak q} \in \text{Ass}_{\overline{S}}(\overline{N})
& \Leftrightarrow
\overline{\mathfrak q}' \in \text{WeakAss}_{\overline{S}'}(\overline{N}')
\end{align*}
Our careful choice of $g$ and the formula for
$\text{Ass}_{\overline{S}'}(\overline{N}')$ above shows that
\begin{equation}
\label{equation-key-observation}
\text{if }\mathfrak r' \in \text{Ass}_{\overline{S}'}(\overline{N}')
\text{ lies over }\mathfrak r \subset \overline{S}\text{ then }
\mathfrak r \subset \overline{\mathfrak q}
\end{equation}
This will be a key observation later in the proof. We will use
the characterization of weakly associated primes given in
Algebra, Lemma \ref{algebra-lemma-weakly-ass-local} without further mention.

\medskip\noindent
Suppose that
$\overline{\mathfrak q} \not \in \text{Ass}_{\overline{S}}(\overline{N})$.
Then
$\overline{\mathfrak q}' \not \in \text{Ass}_{\overline{S}'}(\overline{N}')$.
By
Algebra, Lemmas \ref{algebra-lemma-ass-zero-divisors},
\ref{algebra-lemma-finite-ass}, and
\ref{algebra-lemma-silly}
there exists an element $\overline{a}' \in \overline{\mathfrak q}'$
which is not a zerodivisor on $\overline{N}'$.
After replacing $\overline{a}'$ by $\lambda \overline{a}'$ for some nonzero
$\lambda \in \kappa(\mathfrak p)$ we can find
$a' \in \mathfrak q'$ mapping to $\overline{a}'$. By
Lemma \ref{lemma-invert-universally-injective}
the map $a' : N'_{\mathfrak p'} \to N'_{\mathfrak p'}$ is
$R'_{\mathfrak p'}$-universally injective. In particular
we see that $a' : M' \otimes_{R'} N' \to M' \otimes_{R'} N'$ is
injective after localizing at $\mathfrak p'$ and hence after
localizing at $\mathfrak q'$. Clearly this implies that
$\mathfrak q' \not \in \text{WeakAss}_{S'}(M' \otimes_{R'} N')$.
We conclude that $\mathfrak q \in \text{WeakAss}_S(M \otimes_R N)$ implies
$\overline{\mathfrak q} \in \text{Ass}_{\overline{S}}(\overline{N})$.

\medskip\noindent
Assume $\mathfrak q \in \text{WeakAss}_S(M \otimes_R N)$. We want
to show $\mathfrak p \in \text{WeakAss}_S(M)$.
Let $z \in M \otimes_R N$ be an element such that $\mathfrak q$
is minimal over $J = \text{Ann}_S(z)$.
Let $f_i \in \mathfrak p$, $i \in I$ be a set of generators of the
ideal $\mathfrak p$. Since $\mathfrak q$ lies over $\mathfrak p$, for every $i$
we can choose an $n_i \geq 1$ and $g_i \in S$, $g_i \not \in \mathfrak q$
with $g_i f_i^{n_i} \in J$, i.e., $g_i f_i^{n_i} z = 0$.
Let $z' \in (M' \otimes_{R'} N')_{\mathfrak p'}$ be the image of $z$.
Observe that $z'$ is nonzero because $z$ has nonzero image in
$(M \otimes_R N)_\mathfrak q$ and because $S_\mathfrak q \to S'_{\mathfrak q'}$
is faithfully flat. We claim that $f_i^{n_i} z' = 0$.

\medskip\noindent
Proof of the claim: Let $g'_i \in S'$ be the image of $g_i$.
By the key observation (\ref{equation-key-observation})
we find that the image $\overline{g}'_i \in \overline{S}'$
is not contained in $\mathfrak r'$ for any
$\mathfrak r' \in \text{Ass}_{\overline{S}'}(\overline{N})$.
Hence by Lemma \ref{lemma-invert-universally-injective}
we see that $g'_i : N'_{\mathfrak p'} \to N'_{\mathfrak p'}$ is
$R'_{\mathfrak p'}$-universally injective. In particular
we see that $g'_i : M' \otimes_{R'} N' \to M' \otimes_{R'} N'$ is
injective after localizating at $\mathfrak p'$. The claim
follows because $g_i f_i^{n_i} z' = 0$.

\medskip\noindent
Our claim shows that the annihilator of $z'$ in $R'_{\mathfrak p'}$
contains the elements $f_i^{n_i}$. As $R \to R'$ is \'etale we have
$\mathfrak p'R'_{\mathfrak p'} = \mathfrak pR'_{\mathfrak p'}$
by Algebra, Lemma \ref{algebra-lemma-etale-at-prime}.
Hence the annihilator of $z'$ in $R'_{\mathfrak p'}$ has radical equal to
$\mathfrak p' R_{\mathfrak p'}$ (here we use $z'$ is not zero).
On the other hand
$$
z' \in (M' \otimes_{R'} N')_{\mathfrak p'} =
M'_{\mathfrak p'} \otimes_{R'_{\mathfrak p'}} N'_{\mathfrak p'}
$$
The module $N'_{\mathfrak p'}$ is projective over the local ring
$R'_{\mathfrak p'}$ and hence free
(Algebra, Theorem \ref{algebra-theorem-projective-free-over-local-ring}).
Thus we can find a finite free direct summand $F' \subset N'_{\mathfrak p'}$
such that $z' \in M'_{\mathfrak p'} \otimes_{R'_{\mathfrak p'}} F'$.
If $F'$ has rank $n$, then we deduce that
$\mathfrak p' R'_{\mathfrak p'} \in
\text{WeakAss}_{R'_{\mathfrak p'}}({M'_{\mathfrak p'}}^{\oplus n})$.
This implies
$\mathfrak p'R'_{\mathfrak p'} \in \text{WeakAss}(M'_{\mathfrak p'})$
for example by Algebra, Lemma \ref{algebra-lemma-weakly-ass}.
Then $\mathfrak p' \in \text{WeakAss}_{R'}(M')$
which in turn gives $\mathfrak p \in \text{WeakAss}_R(M)$.
This finishes the proof of the implication
``$\Rightarrow$'' of the equivalence of the lemma.

\medskip\noindent
Assume that $\mathfrak p \in \text{WeakAss}_R(M)$ and
$\overline{\mathfrak q} \in \text{Ass}_{\overline{S}}(\overline{N})$.
We want to show that $\mathfrak q$ is weakly associated to $M \otimes_R N$.
Note that $\overline{\mathfrak q}'$ is a maximal element of
$\text{Ass}_{\overline{S}'}(\overline{N}')$.
This is a consequence of (\ref{equation-key-observation})
and the fact that there are no inclusions among the primes
of $\overline{S}'$ lying over $\overline{\mathfrak q}$
(as fibres of \'etale morphisms are discrete
Morphisms, Lemma \ref{morphisms-lemma-etale-over-field}).
Thus, after replacing $R, S, \mathfrak p, \mathfrak q, M, N$ by
$R', S', \mathfrak p', \mathfrak q', M', N'$
we may assume, in addition to the assumptions of the lemma, that
\begin{enumerate}
\item $\mathfrak p \in \text{WeakAss}_R(M)$,
\item $\overline{\mathfrak q} \in \text{Ass}_{\overline{S}}(\overline{N})$,
\item $N$ is projective as an $R$-module, and
\item $\overline{\mathfrak q}$ is maximal in
$\text{Ass}_{\overline{S}}(\overline{N})$.
\end{enumerate}
There is one more reduction, namely, we may replace
$R, S, M, N$ by their localizations at $\mathfrak p$.
This leads to one more condition, namely,
\begin{enumerate}
\item[(5)] $R$ is a local ring with maximal ideal $\mathfrak p$.
\end{enumerate}
We will finish by showing that (1) -- (5) imply
$\mathfrak q \in \text{WeakAss}(M \otimes_R N)$.

\medskip\noindent
Since $R$ is local and $\mathfrak p \in \text{WeakAss}_R(M)$
we can pick a $y \in M$ whose annihilator $I$ has radical
equal to $\mathfrak p$.
Write $\overline{\mathfrak q} = (\overline{g}_1, \ldots, \overline{g}_n)$
for some $\overline{g}_i \in \overline{S}$. Choose $g_i \in S$
mapping to $\overline{g}_i$.
Then $\mathfrak q = \mathfrak pS + g_1S + \ldots + g_nS$.
Consider the map
$$
\Psi : N/IN \longrightarrow (N/IN)^{\oplus n}, \quad
z \longmapsto (g_1z, \ldots, g_nz).
$$
This is a homomorphism of projective $R/I$-modules.
The local ring $R/I$ is auto-associated
(More on Algebra, Definition \ref{more-algebra-definition-auto-ass})
as $\mathfrak p/I$ is locally nilpotent.
The map $\Psi \otimes \kappa(\mathfrak p)$ is not injective, because
$\overline{\mathfrak q} \in \text{Ass}_{\overline{S}}(\overline{N})$.
Hence More on Algebra, Lemma \ref{more-algebra-lemma-P-fPD-zero}
implies $\Psi$ is not injective. Pick $z \in N/IN$
nonzero in the kernel of $\Psi$. The annihilator $J = \text{Ann}_S(z)$
contains $IS$ and $g_i$ by construction. Thus
$\sqrt{J} \subset S$ contains $\mathfrak q$.
Let $\mathfrak s \subset S$ be a prime minimal over $J$.
Then $\mathfrak q \subset \mathfrak s$,
$\mathfrak s$ lies over $\mathfrak p$, and
$\mathfrak s \in \text{WeakAss}_S(N/IN)$.
The last fact by definition of weakly associated primes.
Apply the ``$\Rightarrow$'' part of the lemma (which we've already proven)
to the ring map $R \to S$ and the modules $R/I$ and $N$
to conclude that
$\overline{\mathfrak s} \in \text{Ass}_{\overline{S}}(\overline{N})$.
Since $\overline{\mathfrak q} \subset \overline{\mathfrak s}$
the maximality of $\overline{\mathfrak q}$, see condition (4) above,
implies that $\overline{\mathfrak q} = \overline{\mathfrak s}$.
This shows that $\mathfrak q = \mathfrak s$ and we conlude
what we want.
\end{proof}

\end{document}
